\section{Thông tin chung}

\subsection{Thông tin nhóm}
\begin{itemize}[leftmargin=1.5cm]
    \item \textbf{Học phần:} Vật lý cho Công nghệ Thông tin
    \item \textbf{Lớp học phần:} \missinginfo{Mã lớp học phần}
    \item \textbf{Giảng viên hướng dẫn:} \missinginfo{ThS. Lê Quốc Hòa (hoặc GVHD phụ trách)}
    \item \textbf{Mã nhóm:} \missinginfo{Mã số nhóm}
\end{itemize}

\vspace{0.3cm}
\textbf{Danh sách thành viên}
\begin{table}[htbp]
\centering
\begin{tabular}{|c|l|l|}
\hline
\textbf{STT} & \textbf{MSSV} & \textbf{Họ và tên} \\ \hline
1 & \missinginfo{MSSV Nguyễn Trần Tuấn Minh} & Nguyễn Trần Tuấn Minh \\ \hline
2 & \missinginfo{MSSV Nguyễn Quốc Dương} & Nguyễn Quốc Dương \\ \hline
3 & \missinginfo{MSSV La Anh Hào} & La Anh Hào \\ \hline
\end{tabular}
\caption{Danh sách thành viên thực hiện đồ án}
\label{tab:thanh_vien}
\end{table}

\subsection{Thông tin đồ án}

\textbf{Tên đề tài:} \textbf{Trạm học tập Pomodoro (Pomodoro Box)}

\vspace{0.2cm}
\textbf{Mô tả sản phẩm:}

Sản phẩm \textbf{Trạm học tập Pomodoro} là một hệ thống thiết bị IoT để bàn kết hợp phần cứng và nền tảng web nhằm tối ưu hóa sự tập trung và quản lý thời gian học tập theo phương pháp Pomodoro. Hệ thống đóng gói trong một hộp vật lý đặt tại góc bàn học, sử dụng vi điều khiển ESP32 kết nối mạng Internet qua Wi-Fi và giao tiếp với hệ thống quản trị trung tâm Node-RED qua giao thức MQTT.

Hệ thống cung cấp giải pháp hỗ trợ học tập toàn diện không phụ thuộc vào điện thoại thông minh:
\begin{itemize}
    \item \textbf{Quản lý phiên học không xao nhãng:} Người dùng kích hoạt phiên học bằng thẻ RFID (RC522) hoặc cài đặt thời lượng phiên học từ xa trên giao diện web. ESP32 thực hiện đếm giờ độc lập và kích hoạt còi (buzzer) khi kết thúc phiên.
    \item \textbf{Trạng thái tập trung trực quan:} Đèn LED đỏ trên mặt hộp được bật tắt bằng nút bấm vật lý để cảnh báo trạng thái ``Đang tập trung'' cho người xung quanh, giảm thiểu gián đoạn.
    \item \textbf{Màn hình tương tác:} Màn hình LCD 1602 I2C hiển thị lời nhắc mục tiêu học tập gửi từ giao diện web cùng thời gian đếm ngược thực tế.
    \item \textbf{Giám sát môi trường học tập:} Cảm biến nhiệt độ, độ ẩm (DHT22) và cảm biến ánh sáng (quang trở LM393) liên tục theo dõi điều kiện phòng học và gửi dữ liệu lên máy chủ.
    \item \textbf{Lưu trữ và phân tích đám mây:} Dữ liệu phiên học và môi trường được lưu trữ theo chuỗi thời gian trên nền tảng Supabase (PostgreSQL), từ đó vẽ biểu đồ trực quan trên Node-RED Dashboard và chạy mô hình phân tích điều kiện học tập tối ưu.
    \item \textbf{Thông báo đa kênh:} Tích hợp gửi tin nhắn cảnh báo qua Telegram Bot khi học quá thời gian quy định mà không nghỉ ngơi, đồng thời gửi báo cáo thống kê qua Gmail.
\end{itemize}

\vspace{0.2cm}
\textbf{Lý do thực hiện:}

Phương pháp Pomodoro đòi hỏi người học chia nhỏ thời gian học thành các phiên 25 phút xen kẽ 5 phút nghỉ ngắn. Tuy nhiên, trong thực tế sinh viên thường gặp nhiều rào cản:
\begin{itemize}
    \item Việc dùng smartphone để bấm giờ thường xuyên gây tác dụng ngược do các thông báo mạng xã hội làm gián đoạn sự chú ý.
    \item Người học có xu hướng quên nghỉ ngơi, dẫn đến căng thẳng thị giác và giảm hiệu suất tiếp thu.
    \item Thiếu sự ghi nhận lịch sử học tập khách quan và không nắm bắt được môi trường học (nhiệt độ quá nóng, phòng thiếu sáng) ảnh hưởng như thế nào đến khả năng tập trung.
    \item Người thân hoặc bạn cùng phòng không biết thời điểm người học đang tập trung cao độ nên vô tình làm phiền.
\end{itemize}
Vì vậy, việc chế tạo một trạm thiết bị vật lý chuyên dụng kết hợp IoT và lưu trữ đám mây giúp giải quyết triệt để các vấn đề trên, tạo ra môi trường học tập kỷ luật, khoa học và tiện lợi.

\subsection{Kiến trúc và luồng truyền nhận dữ liệu}

Hệ thống được thiết kế theo kiến trúc phân tầng chuẩn IoT gồm 4 thành phần chính: Thiết bị phần cứng (ESP32 Box), Trung gian truyền thông (MQTT Broker), Máy chủ điều phối \& Giao diện (Node-RED + Dashboard), và Cơ sở dữ liệu đám mây (Supabase Cloud).

\begin{figure}[htbp]
\centering
\begin{tikzpicture}[
    node distance=1.2cm and 1.5cm,
    box/.style={draw, rounded corners, minimum width=2.2cm, minimum height=1cm, align=center, font=\small},
    hw/.style={box, fill=blue!10, draw=blue!70},
    mid/.style={box, fill=orange!10, draw=orange!70},
    serv/.style={box, fill=green!10, draw=green!70},
    cloud/.style={box, fill=purple!10, draw=purple!70},
    arrow/.style={-{Stealth[scale=1.2]}, thick}
]
    % Node phần cứng
    \node[hw] (esp) {\textbf{ESP32 DevKit}\\(Vi điều khiển)};
    \node[hw, left=of esp, yshift=1.2cm] (sensors) {DHT22, LDR,\\RFID, Nút nhấn\\(INPUT)};
    \node[hw, left=of esp, yshift=-1.2cm] (outputs) {Buzzer, LED đỏ,\\LCD 1602 I2C\\(OUTPUT)};
    
    % Node truyền thông
    \node[mid, right=of esp] (broker) {\textbf{MQTT Broker}\\(HiveMQ / Mosquitto)\\(TCP/1883)};
    
    % Node Server & Dashboard
    \node[serv, right=of broker] (nodered) {\textbf{Node-RED Server}\\(Flow Engine +\\Dashboard Web)};
    
    % Node Cloud & Notify
    \node[cloud, above=of nodered] (supabase) {\textbf{Supabase (Cloud)}\\(PostgreSQL Database)};
    \node[cloud, right=of nodered, yshift=0.9cm] (telegram) {\textbf{Telegram API}\\(Cảnh báo tức thì)};
    \node[cloud, right=of nodered, yshift=-0.9cm] (email) {\textbf{Gmail SMTP}\\(Báo cáo tuần)};

    % Đường kết nối
    \draw[arrow] (sensors) -- (esp);
    \draw[arrow] (esp) -- (outputs);
    \draw[arrow, <->] (esp) -- node[above, font=\footnotesize] {MQTT/Wi-Fi} (broker);
    \draw[arrow, <->] (broker) -- node[above, font=\footnotesize] {Sub/Pub} (nodered);
    \draw[arrow, <->] (nodered) -- node[right, font=\footnotesize] {REST / SQL} (supabase);
    \draw[arrow] (nodered) -- (telegram);
    \draw[arrow] (nodered) -- (email);
\end{tikzpicture}
\caption{Sơ đồ kiến trúc và luồng truyền nhận dữ liệu của hệ thống}
\label{fig:kien_truc_tong_the}
\end{figure}

\vspace{0.2cm}
\textbf{Quy trình hoạt động tổng thể:}
\begin{enumerate}
    \item Khi bật nguồn, ESP32 đọc cấu hình Wi-Fi từ bộ nhớ NVS. Nếu chưa có kết nối, ESP32 kích hoạt chế độ AP phát Wi-Fi cho người dùng cài đặt.
    \item Sau khi kết nối Wi-Fi thành công, ESP32 kết nối đến MQTT Broker và subscribe các topic điều khiển (\texttt{study/timer}, \texttt{study/lcd}).
    \item Định kỳ (mỗi 10 giây), ESP32 đọc cảm biến DHT22 và quang trở LM393, đóng gói định dạng JSON và publish lên topic \texttt{study/dht} và \texttt{study/light}.
    \item Khi người học quẹt thẻ RFID RC522, ESP32 đọc mã UID và publish lên topic \texttt{study/rfid} để thông báo bắt đầu hoặc kết thúc phiên học.
    \item Người học có thể nhấn nút vật lý trên hộp bất kỳ lúc nào để kích hoạt chế độ tập trung (bật sáng đèn LED đỏ độc lập).
    \item Máy chủ Node-RED đón nhận các thông điệp MQTT, giải mã JSON và hiển thị dữ liệu thời gian thực lên giao diện Node-RED Dashboard.
    \item Node-RED thực hiện ghi nhận các chuỗi dữ liệu môi trường và lịch sử phiên học vào các bảng tương ứng trên Supabase.
    \item Từ giao diện Dashboard, người dùng có thể gửi lệnh đếm ngược thời gian xuống ESP32 để kích hoạt còi báo, hoặc gửi lời nhắc hiển thị lên màn hình LCD.
    \item Nếu phát hiện phiên học kéo dài vượt ngưỡng an toàn mà chưa nghỉ, Node-RED tự động kích hoạt bot Telegram gửi tin nhắn cảnh báo.
\end{enumerate}

\subsection{Giải thích flow Website và các node chính}

Hệ thống phần mềm trung tâm vận hành trên nền tảng Node-RED với các khối chức năng cụ thể:
\begin{itemize}
    \item \textbf{Node-RED Dashboard:} Giao diện người dùng tương tác trên nền trình duyệt web, cung cấp các widget hiển thị đồ thị nhiệt độ, độ ẩm, gauge đo cường độ sáng, bảng lịch sử phiên học, ô nhập lời nhắc LCD và thanh điều khiển thời gian Pomodoro.
    \item \textbf{MQTT In / Out Nodes:} Đóng vai trò cầu nối giao tiếp với MQTT Broker. Đảm bảo tiếp nhận các gói tin cảm biến và phân phối các lệnh điều khiển từ web xuống phần cứng với chất lượng dịch vụ QoS 1.
    \item \textbf{Function Nodes:} Thực hiện tiền xử lý dữ liệu, chuẩn hóa cấu trúc JSON, tính toán mức sáng định tính (``du\_sang'' / ``thieu\_sang''), và quản lý trạng thái phiên học.
    \item \textbf{Supabase Integration Nodes:} Kết nối với dịch vụ cơ sở dữ liệu Supabase thông qua REST API hoặc truy vấn trực tiếp PostgreSQL để thực hiện các thao tác thêm mới và truy xuất lịch sử.
    \item \textbf{Notification Nodes:} Sử dụng node Telegram và Email để kết nối với dịch vụ của bên thứ ba, thực hiện cảnh báo tự động khi thỏa mãn điều kiện logic được định sẵn.
\end{itemize}

\subsection{Thiết kế cơ sở dữ liệu}

Hệ thống sử dụng cơ sở dữ liệu quan hệ PostgreSQL trên nền tảng đám mây Supabase với ba bảng chính:

\begin{table}[htbp]
\centering
\begin{tabular}{|l|p{11cm}|}
\hline
\textbf{Tên bảng} & \textbf{Nội dung lưu trữ chính} \\ \hline
\texttt{users} & ID người dùng, tên đăng nhập, mã băm mật khẩu (bcrypt), mã thẻ RFID UID, thông tin liên hệ (email, chat ID Telegram) và thời điểm khởi tạo. \\ \hline
\texttt{study\_sessions} & ID phiên học, liên kết user, thời điểm bắt đầu, thời điểm kết thúc, thời lượng phiên học (phút), trạng thái hoàn thành (completed / interrupted). \\ \hline
\texttt{sensor\_logs} & ID bản ghi, thời điểm ghi nhận (timestamp), nhiệt độ, độ ẩm môi trường từ DHT22, giá trị thô và phân loại ánh sáng từ quang trở. \\ \hline
\end{tabular}
\caption{Các bảng dữ liệu chính trên Supabase}
\label{tab:database_tables}
\end{table}

\begin{figure}[htbp]
\centering
\begin{tikzpicture}[
    node distance=1.5cm and 2.5cm,
    entity/.style={draw, rectangle, minimum width=3.5cm, minimum height=2.2cm, align=left, font=\footnotesize, fill=gray!5},
    title/.style={font=\bfseries\small, align=center}
]
    \node[entity] (users) {
        \textbf{USERS}\\[2pt]
        \rule{3.2cm}{0.4pt}\\[2pt]
        \underline{id}: uuid [PK]\\
        username: text\\
        password\_hash: text\\
        rfid\_uid: text\\
        telegram\_chat\_id: text\\
        created\_at: timestamptz
    };
    
    \node[entity, right=of users] (sessions) {
        \textbf{STUDY\_SESSIONS}\\[2pt]
        \rule{3.2cm}{0.4pt}\\[2pt]
        \underline{id}: uuid [PK]\\
        user\_id: uuid [FK]\\
        start\_time: timestamptz\\
        end\_time: timestamptz\\
        duration\_minutes: int\\
        status: text
    };
    
    \node[entity, below=of users] (logs) {
        \textbf{SENSOR\_LOGS}\\[2pt]
        \rule{3.2cm}{0.4pt}\\[2pt]
        \underline{id}: uuid [PK]\\
        logged\_at: timestamptz\\
        temperature: float\\
        humidity: float\\
        light\_raw: int\\
        light\_level: text
    };

    \draw[thick] (users.east) -- node[above, font=\scriptsize] {1} node[below, font=\scriptsize] {quản lý} (sessions.west) node[above, font=\scriptsize] {N};
\end{tikzpicture}
\caption{Sơ đồ cấu trúc và quan hệ cơ sở dữ liệu trên Supabase}
\label{fig:erd_database}
\end{figure}

\subsection{Đối chiếu yêu cầu tối thiểu của đồ án}

Dự án đáp ứng đầy đủ và vượt mức 6 yêu cầu bắt buộc của học phần Vật lý cho Công nghệ Thông tin:
\begin{enumerate}
    \item \textbf{Bo mạch có khả năng kết nối Internet:} Hệ thống sử dụng vi điều khiển ESP32 DevKit V1 tích hợp sẵn Wi-Fi 2.4 GHz, trao đổi dữ liệu hai chiều thông qua giao thức MQTT qua Internet.
    \item \textbf{Có ít nhất một thiết bị INPUT và một thiết bị OUTPUT:}
    \begin{itemize}
        \item \textit{Thiết bị INPUT:} Cảm biến nhiệt độ - độ ẩm DHT22, Module quang trở LM393, Đầu đọc RFID RC522, Nút nhấn cơ học.
        \item \textit{Thiết bị OUTPUT:} Còi báo (Buzzer chủ động), Đèn LED đỏ báo trạng thái tập trung, Màn hình LCD 1602 I2C.
        \item Tất cả các thiết bị đôi một khác nhau, không xảy ra trùng lặp.
    \end{itemize}
    \item \textbf{Hiển thị thông tin của thiết bị đầu vào trên Website qua Internet:} Giao diện Node-RED Dashboard hiển thị trực quan thông số nhiệt độ, độ ẩm từ DHT22 và mức sáng phòng học từ quang trở theo thời gian thực; hiển thị mã UID và trạng thái phiên học khi quét thẻ RFID.
    \item \textbf{Điều khiển thiết bị đầu ra từ Website qua Internet:} Người dùng thiết lập số phút hẹn giờ Pomodoro trên web để gửi lệnh xuống ESP32 kích hoạt còi buzzer khi hết giờ; gửi nội dung văn bản tùy chỉnh lên màn hình LCD 1602.
    \item \textbf{Lưu trữ dữ liệu trên Cloud:} Toàn bộ thông số môi trường theo thời gian và lịch sử phiên học được lưu trữ tập trung trên hệ thống đám mây Supabase PostgreSQL.
    \item \textbf{Thông báo đến điện thoại hoặc Email trong tình huống khẩn cấp:} Tự động kích hoạt bot gửi tin nhắn cảnh báo qua Telegram khi người học tập trung liên tục vượt quá 90 phút mà không giải lao; định kỳ gửi email tổng kết hiệu quả học tập mỗi tuần qua Gmail SMTP.
\end{enumerate}

\subsection{Danh sách chức năng}

\begin{longtable}{|c|p{4.5cm}|p{8.5cm}|}
\hline
\textbf{Mã} & \textbf{Chức năng} & \textbf{Mô tả chi tiết và Luồng xử lý} \\ \hline
\endhead

\textbf{CN01} & Giám sát nhiệt độ \& độ ẩm & Đọc dữ liệu từ DHT22 gửi lên Node-RED qua topic \texttt{study/dht} và hiển thị số đo góc học tập trên Dashboard. \\ \hline

\textbf{CN02} & Giám sát mức độ ánh sáng & Đọc giá trị tương tự từ quang trở LM393 gửi lên \texttt{study/light}, phân loại đủ sáng / thiếu sáng. \\ \hline

\textbf{CN03} & Hẹn giờ Pomodoro với còi báo & Dashboard gửi số phút lên topic \texttt{study/timer}, ESP32 đếm giờ độc lập và kích hoạt còi buzzer khi kết thúc phiên. \\ \hline

\textbf{CN04} & Bật/tắt đèn báo tập trung & Nhấn nút cơ học trên hộp để bật/tắt đèn LED đỏ báo hiệu trạng thái tập trung cho người xung quanh (xử lý trực tiếp trên ESP32). \\ \hline

\textbf{CN05} & Quản lý phiên học bằng RFID & Quẹt thẻ RFID RC522 gửi UID lên \texttt{study/rfid} để bắt đầu hoặc kết thúc phiên Pomodoro, cập nhật trạng thái lên Dashboard. \\ \hline

\textbf{CN06} & Hiển thị lời nhắc trên LCD & Dashboard gửi chuỗi ký tự lời nhắc/mục tiêu lên topic \texttt{study/lcd}, ESP32 hiển thị lên màn hình LCD 1602 I2C. \\ \hline

\textbf{CN07} & Lưu trữ dữ liệu đám mây & Node-RED tự động ghi chuỗi thời gian cảm biến và lịch sử phiên học vào cơ sở dữ liệu Supabase PostgreSQL. \\ \hline

\textbf{CN08} & Biểu đồ thống kê lịch sử & Dashboard đọc dữ liệu lịch sử từ Supabase để vẽ biểu đồ nhiệt độ, độ ẩm và thống kê số phiên Pomodoro theo ngày/tuần. \\ \hline

\textbf{CN09} & Cảnh báo học quá lâu qua Telegram & Tự động gửi tin nhắn Telegram khi một phiên học vượt quá 90 phút nhằm nhắc nhở giải lao bảo vệ sức khỏe. \\ \hline

\textbf{CN10} & Báo cáo tuần qua Gmail & Node-RED tổng hợp số phiên Pomodoro và thời gian học tập trong tuần, tự động gửi báo cáo qua Gmail SMTP. \\ \hline

\textbf{CN11} & Xác thực người dùng Dashboard & Hệ thống đăng ký, đăng nhập tài khoản bảo mật bằng bcrypt lưu trên Supabase; phân quyền xem và điều khiển hệ thống. \\ \hline

\textbf{CN12} & Mô hình phân tích điều kiện học & Thuật toán tự xây dựng phân tích mối liên hệ giữa nhiệt độ, ánh sáng, khung giờ với tỷ lệ hoàn thành phiên học. \\ \hline

\textbf{CN13} & Cấu hình Wi-Fi qua Web Server & Khi mất Wi-Fi, ESP32 phát mạng AP \texttt{PomodoroBox-Setup} và mở web server nội bộ \texttt{192.168.4.1} cho người dùng cài đặt. \\ \hline
\caption{Danh sách các chức năng chính của hệ thống Trạm học tập Pomodoro}
\label{tab:danh_sach_chuc_nang}
\end{longtable}
